\documentclass[aspectratio=169,14pt]{beamer}

% Compile with XeLaTeX, LuaLaTeX, or Tectonic.
\usepackage[UTF8,scheme=plain,fontset=fandol]{ctex}
\usepackage{amsmath,amssymb}
\usepackage{booktabs,tabularx,array,colortbl}

\usetheme{Madrid}
\usefonttheme{professionalfonts}

\definecolor{BeamerRoyal}{HTML}{3939B8}
\definecolor{BeamerDeep}{HTML}{171B72}
\definecolor{BeamerInk}{HTML}{20212B}
\definecolor{SoftBlue}{HTML}{F1F2FF}
\definecolor{SoftGray}{HTML}{F5F5F7}
\definecolor{WarmAccent}{HTML}{D85A3A}

\setbeamercolor{normal text}{fg=BeamerInk,bg=white}
\setbeamercolor{structure}{fg=BeamerRoyal}
\setbeamercolor{frametitle}{fg=white,bg=BeamerRoyal}
\setbeamercolor{title}{fg=white,bg=BeamerRoyal}
\setbeamercolor{subtitle}{fg=BeamerInk}
\setbeamercolor{block title}{fg=white,bg=BeamerRoyal}
\setbeamercolor{block body}{fg=BeamerInk,bg=SoftBlue}
\setbeamercolor{alerted text}{fg=WarmAccent}
\setbeamercolor{author in head/foot}{fg=white,bg=BeamerDeep}
\setbeamercolor{title in head/foot}{fg=white,bg=BeamerRoyal}
\setbeamercolor{date in head/foot}{fg=white,bg=BeamerDeep}

\setbeamertemplate{items}[ball]
\setbeamertemplate{blocks}[rounded][shadow=true]
\setbeamertemplate{navigation symbols}{}
\setbeamertemplate{headline}{}
\setbeamertemplate{title page}{%
  \vspace{0.55cm}
  \begin{beamercolorbox}[wd=.92\paperwidth,sep=.34cm,center,rounded=true,shadow=true]{title}
    \usebeamerfont{title}\inserttitle
  \end{beamercolorbox}
  \vspace{0.38cm}
  \begin{center}
    {\usebeamerfont{subtitle}\insertsubtitle\par}
    \vspace{0.70cm}
    {\usebeamerfont{author}\insertauthor\par}
    \vspace{0.18cm}
    {\usebeamerfont{institute}\insertinstitute\par}
    \vspace{0.52cm}
    {\usebeamerfont{date}\insertdate\par}
  \end{center}
}

\setbeamertemplate{frametitle}{%
  \nointerlineskip
  \begin{beamercolorbox}[wd=\paperwidth,ht=.72cm,dp=.18cm,leftskip=.38cm]{frametitle}
    \usebeamerfont{frametitle}\insertframetitle
  \end{beamercolorbox}
}

\setbeamertemplate{footline}{%
  \leavevmode
  \hbox{%
    \begin{beamercolorbox}[wd=.31\paperwidth,ht=2.6ex,dp=1.05ex,center]{author in head/foot}
      \usebeamerfont{author in head/foot}\insertshortauthor\enspace(\insertshortinstitute)
    \end{beamercolorbox}%
    \begin{beamercolorbox}[wd=.36\paperwidth,ht=2.6ex,dp=1.05ex,center]{title in head/foot}
      \usebeamerfont{title in head/foot}\insertshorttitle
    \end{beamercolorbox}%
    \begin{beamercolorbox}[wd=.33\paperwidth,ht=2.6ex,dp=1.05ex,center]{date in head/foot}
      \usebeamerfont{date in head/foot}\insertshortdate\hspace{1.2em}
      \insertframenumber/\inserttotalframenumber
    \end{beamercolorbox}%
  }
}

\setbeamersize{text margin left=.68cm,text margin right=.68cm}
\setbeamerfont{title}{size*={29pt}{35pt},series=\bfseries}
\setbeamerfont{subtitle}{size*={17pt}{22pt}}
\setbeamerfont{author}{size*={16pt}{20pt},series=\bfseries}
\setbeamerfont{institute}{size*={13pt}{17pt}}
\setbeamerfont{date}{size*={14pt}{18pt}}
\setbeamerfont{frametitle}{size*={22pt}{26pt}}
\setbeamerfont{normal text}{size*={16pt}{20pt}}
\setbeamerfont{itemize/enumerate body}{size*={15.5pt}{20pt}}
\setbeamerfont{itemize/enumerate subbody}{size*={14pt}{18pt}}
\setbeamerfont{block title}{size*={15pt}{19pt},series=\bfseries}
\setbeamerfont{block body}{size*={14.5pt}{19pt}}
\setbeamerfont{author in head/foot}{size*={7.5pt}{9pt}}
\setbeamerfont{title in head/foot}{size*={7.5pt}{9pt}}
\setbeamerfont{date in head/foot}{size*={7.5pt}{9pt}}

\newcommand{\key}[1]{\textcolor{BeamerRoyal}{\textbf{#1}}}
\newcommand{\unit}[1]{\,\mathrm{#1}}
\newcommand{\takeaway}[1]{%
  \begin{beamercolorbox}[wd=\linewidth,sep=.22cm,center,rounded=true,shadow=true]{block body}
    \textcolor{BeamerDeep}{\textbf{#1}}
  \end{beamercolorbox}
}

\title[Voxel Spacing]{Voxel Spacing 与 3D CNN}
\subtitle{医学影像预处理中的统一物理尺度}
\author[学习笔记]{NV-Segment-CT/MRI 学习笔记}
\institute[医学影像分割]{面向入门学习者的简明讲解}
\date[2026.09]{2026 年 9 月}

\begin{document}

\begin{frame}
  \titlepage
\end{frame}

\begin{frame}{内容概览}
  \vspace{0.15cm}
  \begin{enumerate}
    \item \key{定义}：一个 voxel 在现实中对应多大距离
    \item \key{区别}：voxel 数量和物理尺寸是两回事
    \item \key{预处理}：resampling 与 interpolation 如何统一网格
    \item \key{模型影响}：spacing 为什么会改变 3D CNN 看到的形状
  \end{enumerate}
  \vfill
  \takeaway{主线：从“数据格子”回到“真实毫米”}
\end{frame}

\begin{frame}{Voxel Spacing 的含义}
  \begin{columns}[T,onlytextwidth]
    \column{.54\textwidth}
      \begin{itemize}
        \item spacing 记录相邻 voxel 在 $x,y,z$ 三个方向对应的\key{真实距离}
        \item 单位通常是 $\mathrm{mm/voxel}$
        \item 三个方向数值不同，voxel 可近似理解为长方体
      \end{itemize}
    \column{.41\textwidth}
      \begin{block}{例：一份 CT}
        \centering
        \vspace{0.08cm}
        $\mathbf{s}=(0.7,\ 0.7,\ 5.0)\unit{mm}$
        \vspace{0.12cm}

        \renewcommand{\arraystretch}{1.22}
        \begin{tabular}{c@{\quad}c}
          $x$ 方向 & $0.7\unit{mm}$\\
          $y$ 方向 & $0.7\unit{mm}$\\
          $z$ 方向 & $5.0\unit{mm}$
        \end{tabular}
      \end{block}
  \end{columns}
  \vspace{0.24cm}
  \takeaway{入门记忆：spacing 决定一个 voxel 在现实中有多大}
  \vspace{0.08cm}
  {\fontsize{9.5}{12}\selectfont\color{black!58}
  严格地说，spacing 描述相邻采样点的物理间距。这里把它理解成 voxel 的三边长度。}
\end{frame}

\begin{frame}{矩阵大小与物理范围}
  \begin{columns}[T,onlytextwidth]
    \column{.47\textwidth}
      \begin{block}{矩阵大小 / voxel 数量}
        \centering
        $512\times512\times100$ voxel\\[-0.02cm]
        {\small 一共有多少个数据格子}
      \end{block}
    \column{.47\textwidth}
      \begin{block}{Voxel spacing / 物理尺度}
        \centering
        $0.8\times0.8\times5.0\unit{mm}$\\[-0.02cm]
        {\small 每个格子对应多少毫米}
      \end{block}
  \end{columns}

  \vspace{-0.08cm}
  \[
    \boxed{\mathbf{L}\approx\mathbf{N}\odot\mathbf{s}}
    \qquad
    \text{物理范围}\approx\text{voxel 数量}\times\text{spacing}
  \]
  \vspace{-0.28cm}
  \[
    L_x=512\times0.8=409.6\unit{mm}
    \qquad
    L_z=100\times5=500\unit{mm}
  \]
  \vspace{-0.12cm}
  \takeaway{{\fontsize{14}{18}\selectfont
    只有把矩阵大小和 spacing 放在一起，\\
    才能知道扫描覆盖多大范围}}
\end{frame}

\begin{frame}{扫描参数带来的尺度差异}
  \centering
  {\small
    \renewcommand{\arraystretch}{1.27}
    \begin{tabular}{>{\bfseries}lcc}
      \toprule
        & \textcolor{BeamerRoyal}{CT A} & \textcolor{BeamerRoyal}{CT B}\\
      \midrule
      Spacing & $0.7\times0.7\times5.0\unit{mm}$ & $1.0\times1.0\times1.0\unit{mm}$\\
      $z$ 方向 $10\unit{mm}$ 的结构 & $10/5=2$ voxel & $10/1=10$ voxel\\
      模型看到的厚度 & 约 $2$ 层 & 约 $10$ 层\\
      \bottomrule
    \end{tabular}
  }

  \vspace{0.42cm}
  \takeaway{{\fontsize{14}{18}\selectfont
    同一个真实结构可能显得厚度不同。\\
    模型会把扫描尺度误当成形状差异。}}
\end{frame}

\begin{frame}{Resampling 统一物理网格}
  \begin{columns}[T,onlytextwidth]
    \column{.51\textwidth}
      \begin{block}{原始 $z$ 方向}
        \centering
        $100\text{ voxel}\times5\unit{mm/voxel}=500\unit{mm}$
      \end{block}
      \begin{block}{目标 spacing：$1\unit{mm/voxel}$}
        \centering
        $N_{\mathrm{new}}\approx\dfrac{500\unit{mm}}{1\unit{mm/voxel}}
        =500\text{ voxel}$
      \end{block}
    \column{.44\textwidth}
      \textbf{系统依次完成：}
      \begin{enumerate}
        \item 建立目标物理网格
        \item 找到新 voxel 在原图中的位置
        \item 用插值估计该位置的数值
      \end{enumerate}
  \end{columns}

  \vspace{0.20cm}
  \takeaway{身体覆盖范围仍约为 $500\unit{mm}$，只是表示它的 voxel 更密集}
\end{frame}

\begin{frame}{Interpolation 估计中间值}
  \[
    I(z)=(1-\alpha)I_0+\alpha I_1,
    \qquad \alpha=\frac{z-z_0}{z_1-z_0}
  \]

  \centering
  \renewcommand{\arraystretch}{1.42}
  \begin{tabular}{>{\columncolor{SoftGray}}c*{6}{c}}
    位置 $z$ & $0$ & $1$ & $2$ & $3$ & $4$ & $5\unit{mm}$\\
    \hline
    灰度 $I$ & \cellcolor{BeamerRoyal!18}$100$ & $120$ & $140$ & $160$ & $180$ & \cellcolor{BeamerRoyal!18}$200$\\
  \end{tabular}

  \vspace{0.34cm}
  \begin{columns}[T,onlytextwidth]
    \column{.47\textwidth}
      \begin{block}{CT / MRI 灰度}
        常用线性或三线性插值，让强度平滑变化。
      \end{block}
    \column{.47\textwidth}
      \begin{block}{器官标签 mask}
        常用最近邻插值，避免产生不存在的类别编号。
      \end{block}
  \end{columns}

  \vfill
  \takeaway{新增 voxel 的数值来自估计，不是扫描仪新获得的信息}
\end{frame}

\begin{frame}{Spacing 对 3D CNN 的影响}
  \textbf{同一个 $3\times3\times3$ 卷积核，在不同 spacing 下覆盖的现实范围不同：}

  \vspace{0.24cm}
  \centering
  \renewcommand{\arraystretch}{1.38}
  \begin{tabular}{lcc}
    \toprule
      & Spacing & 卷积核约覆盖的物理范围\\
    \midrule
    CT A & $0.7\times0.7\times5.0\unit{mm}$ & $2.1\times2.1\times15\unit{mm}$\\
    CT B & $1.0\times1.0\times1.0\unit{mm}$ & $3\times3\times3\unit{mm}$\\
    \bottomrule
  \end{tabular}

  \vspace{0.38cm}
  \begin{itemize}
    \item 未统一时，模型会在两份数据中观察不同大小的现实区域
    \item 统一 spacing 后，相同器官更接近相同的 voxel 尺度
    \item 模型因此更容易学习真实形状，减少扫描参数的干扰
  \end{itemize}
\end{frame}

\begin{frame}{信息边界与总结}
  \begin{alertblock}{Resampling 的边界}
    \centering
    原始 $z$ spacing 从 $5\unit{mm}$ 表示为 $1\unit{mm}$，会得到更多切片，\textbf{不会增加真实扫描信息}。
  \end{alertblock}

  \vspace{0.22cm}
  \begin{itemize}
    \item 插值让坐标网格更统一，原始分辨率仍由扫描决定
    \item 原始 $z$ spacing 很大时，分割表面仍可能出现台阶感
    \item spacing 与矩阵大小共同决定影像覆盖的物理空间
  \end{itemize}

  \vfill
  \takeaway{Voxel spacing 决定现实尺寸。Resampling 把不同影像放到统一的物理网格上。}
\end{frame}

\end{document}
